\chapter{Cadrage conceptuel et état de l'art}

\section*{Introduction}
\addcontentsline{toc}{section}{Introduction}

La Migration Factory ne repose pas uniquement sur des outils de mise à niveau. Elle combine des agents spécialisés, un orchestrateur, des services déterministes, des outils externes et une supervision humaine. La compréhension de la solution exige donc un cadrage conceptuel portant à la fois sur les agents intelligents, les systèmes multi-agents, l'orchestration et la migration logicielle.

Ce chapitre présente d'abord les notions fondamentales sous une forme proche d'un cours : définition et propriétés d'un agent, organisation d'un système multi-agents, modes de coordination et modèles d'orchestration. Il expose ensuite les architectures techniques permettant d'intégrer les agents dans une plateforme logicielle. La seconde partie rappelle les concepts de migration et analyse plusieurs travaux, outils et plateformes proches afin de positionner précisément la contribution du projet.

% ============================================================
\section{Agents intelligents}
% ============================================================

\subsection{Définition}

Un agent logiciel est une entité informatique située dans un environnement, capable de percevoir des informations, de maintenir un état ou un contexte, puis d'agir afin de poursuivre un objectif. Wooldridge et Jennings distinguent la théorie des agents, leurs architectures et les langages permettant de les construire \cite{wooldridge-jennings}. Dans une perspective multi-agents, l'agent ne se définit pas seulement par un appel à un modèle de langage : il doit posséder une responsabilité, des capacités, des entrées, des sorties et des règles d'interaction.

\begin{figure}[H]
    \centering
    \resizebox{0.78\textwidth}{!}{%
    \begin{tikzpicture}[
        box/.style={draw=cgiblue,rounded corners=3pt,fill=cgilight!55,minimum width=4.0cm,minimum height=1.2cm,align=center,font=\small\bfseries},
        env/.style={draw=cgigray,rounded corners=3pt,fill=white,minimum width=4.0cm,minimum height=1.2cm,align=center,font=\small\bfseries},
        arr/.style={-{Latex[length=2.2mm]},thick,draw=cgigray}
    ]
        \node[env] (env) {Environnement};
        \node[box,right=2.0cm of env] (agent) {Agent\\état, objectifs, capacités};
        \draw[arr,bend left=18] (env) to node[above,font=\small]{perceptions} (agent);
        \draw[arr,bend left=18] (agent) to node[below,font=\small]{actions} (env);
    \end{tikzpicture}%
    }
    \caption{Interaction générale entre un agent et son environnement}
    \label{fig:agent-environnement-ch2}
\end{figure}

Dans le cadre du projet, les agents sont des composants spécialisés et bornés. Ils ne contrôlent pas librement le système de fichiers ou le terminal. Ils reçoivent un contexte préparé par le backend, produisent un résultat structuré, puis rendent la main à l'orchestrateur.

\subsection{Propriétés d'un agent}

La littérature associe généralement plusieurs propriétés aux agents \cite{wooldridge-book,weiss-mas} :

\begin{itemize}
    \item \textbf{autonomie :} capacité à choisir une action sans intervention continue ;
    \item \textbf{réactivité :} capacité à répondre aux changements de l'environnement ;
    \item \textbf{proactivité :} aptitude à poursuivre un objectif et non à répondre uniquement à un événement ;
    \item \textbf{sociabilité :} capacité à communiquer avec d'autres agents ou acteurs ;
    \item \textbf{raisonnement :} sélection d'une action à partir d'informations et de contraintes ;
    \item \textbf{adaptabilité :} modification du comportement selon le contexte ou les résultats précédents.
\end{itemize}

Ces propriétés doivent être interprétées avec prudence. Tous les agents de la Migration Factory ne sont pas autonomes au même niveau. L'analyseur réagit à une tâche précise ; le planificateur poursuit l'objectif de produire une trajectoire ; le reviewer évalue une sortie ; l'assistant converse avec le développeur. L'autonomie reste limitée par l'orchestrateur, les schémas de sortie et les règles d'autorité.

\subsection{Types d'agents}

\begin{table}[H]
    \centering
    \small
    \arrayrulecolor{cgigray}
    \rowcolors{2}{cgilight!55}{white}
    \begin{tabularx}{\textwidth}{|>{\bfseries\RaggedRight\arraybackslash}p{3.7cm}|>{\RaggedRight\arraybackslash}X|>{\RaggedRight\arraybackslash}p{4.0cm}|}
        \hline
        \rowcolor{cgilight!95}
        \textbf{Type} & \textbf{Rôle général} & \textbf{Exemple dans le projet} \\
        \hline
        Agent réactif & Répond à un événement ou à un état sans construire une stratégie longue. & Agent de qualification ou de contrôle d'un résultat. \\
        \hline
        Agent cognitif & Utilise un modèle du problème et raisonne sur plusieurs informations. & Analyseur et agent de diagnostic. \\
        \hline
        Agent conversationnel & Interagit en langage naturel avec un utilisateur. & Assistant de migration. \\
        \hline
        Agent de planification & Décompose un objectif en tâches ordonnées. & Planificateur de trajectoire. \\
        \hline
        Agent de monitoring & Observe l'état, les événements et les anomalies. & Projection du cockpit et mécanismes de surveillance. \\
        \hline
        Agent d'exécution & Réalise une tâche opérationnelle dans un environnement. & Transformateur borné ; l'exécution shell reste déterministe. \\
        \hline
        Agent évaluateur & Examine la qualité ou la conformité d'une sortie. & Reviewers d'analyse, de plan et de réparation. \\
        \hline
    \end{tabularx}
    \caption{Types d'agents et correspondance avec le projet}
    \label{tab:types-agents-ch2}
    \rowcolors{2}{white}{white}
    \arrayrulecolor{black}
\end{table}

\subsection{Agent, LLM, service et outil}

Un \textbf{LLM} produit une sortie probabiliste à partir d'un contexte. Un \textbf{agent} ajoute à ce modèle un rôle, une mémoire ou un état, des règles d'interaction et éventuellement des outils. Un \textbf{service} implémente une capacité déterministe accessible par contrat. Un \textbf{outil} réalise une opération technique, par exemple Maven, OpenRewrite ou Angular CLI.

Cette distinction est centrale : appeler un LLM ne suffit pas pour construire un agent, et confier une commande à un service déterministe ne transforme pas ce service en agent. La Migration Factory associe ces catégories sans leur attribuer la même autorité.

% ============================================================
\section{Systèmes multi-agents}
% ============================================================

\subsection{Définition}

Un système multi-agents (SMA) regroupe plusieurs agents capables d'interagir dans un environnement partagé afin de poursuivre des objectifs individuels ou collectifs. Les ouvrages de Wooldridge, Weiss et Ferber présentent les SMA comme des systèmes dans lesquels les responsabilités, les connaissances et les actions sont distribuées entre plusieurs entités \cite{wooldridge-book,weiss-mas,ferber-mas}.

\begin{figure}[H]
    \centering
    \resizebox{0.92\textwidth}{!}{%
    \begin{tikzpicture}[
        agent/.style={draw=emsigreen!70!black,rounded corners=3pt,fill=emsigreen!8,minimum width=3.2cm,minimum height=1.05cm,align=center,font=\small\bfseries},
        state/.style={draw=cgiblue,rounded corners=3pt,fill=cgilight!55,minimum width=4.0cm,minimum height=1.05cm,align=center,font=\small\bfseries},
        arr/.style={<->,thick,draw=cgigray}
    ]
        \node[state] (shared) {État et artefacts partagés};
        \node[agent,above left=1.3cm and 1.4cm of shared] (a1) {Agent d'analyse};
        \node[agent,above right=1.3cm and 1.4cm of shared] (a2) {Agent de planification};
        \node[agent,below left=1.3cm and 1.4cm of shared] (a3) {Agent de transformation};
        \node[agent,below right=1.3cm and 1.4cm of shared] (a4) {Agent de révision};
        \draw[arr] (a1) -- (shared);
        \draw[arr] (a2) -- (shared);
        \draw[arr] (a3) -- (shared);
        \draw[arr] (a4) -- (shared);
    \end{tikzpicture}%
    }
    \caption{Vue simplifiée d'un système multi-agents spécialisé}
    \label{fig:sma-simplifie-ch2}
\end{figure}

\subsection{Coordination, communication et délégation}

La \textbf{coordination} organise les dépendances entre actions. La \textbf{communication} transmet des informations, demandes ou décisions. La \textbf{délégation} confie une sous-tâche à un agent disposant des capacités appropriées. La \textbf{coopération} apparaît lorsque plusieurs agents contribuent à un objectif commun. La \textbf{négociation} intervient lorsque plusieurs agents doivent arbitrer des préférences ou ressources.

FIPA a standardisé des structures de messages et des protocoles d'interaction entre agents. Son Agent Communication Language associe à chaque message une intention communicative, telle qu'informer, demander ou proposer \cite{fipa-acl}. La Migration Factory ne met pas en œuvre l'intégralité de FIPA-ACL, mais reprend le principe de contrats explicites : type de tâche, contexte, sortie, statut et corrélation.

\subsection{Avantages et limites}

Une architecture SMA peut apporter :

\begin{itemize}
    \item \textbf{modularité :} chaque agent possède une responsabilité limitée ;
    \item \textbf{scalabilité organisationnelle :} les capacités peuvent être développées séparément ;
    \item \textbf{robustesse :} une défaillance peut être isolée et examinée ;
    \item \textbf{distribution des responsabilités :} proposition, review et exécution sont séparées ;
    \item \textbf{flexibilité :} un agent ou un modèle peut être remplacé sans refonte complète ;
    \item \textbf{gestion du contexte :} chaque agent reçoit seulement les informations utiles.
\end{itemize}

La multiplication des agents introduit néanmoins des coûts : davantage de contrats, d'états, de latence, d'appels LLM et de risques d'incohérence. La documentation LangChain rappelle qu'une tâche complexe ne nécessite pas automatiquement plusieurs agents ; un agent unique bien outillé peut parfois suffire \cite{langchain-multiagent}. Le recours au SMA doit donc être justifié par la spécialisation réelle des rôles et les contraintes de gouvernance.

% ============================================================
\section{Orchestration d'agents}
% ============================================================

\subsection{Orchestration, chorégraphie, coordination et planification}

\begin{table}[H]
    \centering
    \small
    \arrayrulecolor{cgigray}
    \rowcolors{2}{cgilight!55}{white}
    \begin{tabularx}{\textwidth}{|>{\bfseries\RaggedRight\arraybackslash}p{3.1cm}|>{\RaggedRight\arraybackslash}X|>{\RaggedRight\arraybackslash}p{4.1cm}|}
        \hline
        \rowcolor{cgilight!95}
        \textbf{Notion} & \textbf{Définition} & \textbf{Application au projet} \\
        \hline
        Orchestration & Un composant central décide de l'ordre, des transitions et des participants. & Modèle principal de la Migration Factory. \\
        \hline
        Chorégraphie & Les participants réagissent à des messages selon un protocole partagé sans contrôleur unique. & Non retenue comme modèle principal. \\
        \hline
        Coordination & Ensemble des mécanismes évitant les conflits et respectant les dépendances. & États, gates, contrats, priorités et idempotence. \\
        \hline
        Planification & Construction d'une séquence de tâches permettant d'atteindre un objectif. & Trajectoire de versions et plan de migration. \\
        \hline
    \end{tabularx}
    \caption{Distinction entre orchestration, chorégraphie, coordination et planification}
    \label{tab:notions-orchestration-ch2}
    \rowcolors{2}{white}{white}
    \arrayrulecolor{black}
\end{table}

La chorégraphie de services décrit des interactions observables entre participants qui connaissent un protocole commun, sans qu'un seul service pilote toutes les actions \cite{w3c-choreography}. La Migration Factory suit au contraire un état central et des transitions contrôlées ; elle relève donc d'une orchestration centralisée complétée par des agents spécialisés.

\subsection{Modèles d'organisation}

Quatre modèles sont courants :

\begin{itemize}
    \item \textbf{orchestrateur central :} toutes les tâches passent par un coordinateur ;
    \item \textbf{agents distribués :} les agents se coordonnent directement ;
    \item \textbf{architecture hiérarchique :} des superviseurs délèguent à des sous-agents ;
    \item \textbf{architecture hybride :} l'orchestration centralise les états sensibles tandis que les agents conservent une autonomie locale.
\end{itemize}

Le projet adopte le quatrième modèle. L'orchestrateur choisit le prochain nœud selon l'état. L'agent traite sa responsabilité localement. Les services déterministes appliquent les commandes. Le développeur intervient sur les décisions sensibles.

\subsection{Mécanismes d'orchestration}

Les principaux mécanismes sont :

\begin{itemize}
    \item routage d'une tâche vers l'agent approprié ;
    \item séquencement et dépendances entre étapes ;
    \item gestion des priorités et tentatives ;
    \item mémoire ou état partagé ;
    \item supervision et journalisation ;
    \item interruption et reprise ;
    \item gestion des erreurs et compensation ;
    \item replanification à partir d'une nouvelle preuve.
\end{itemize}

LangGraph permet de construire des workflows personnalisés combinant logique déterministe et comportements agentiques, avec un état persistant et des interruptions humaines \cite{langgraph-persistence,langchain-multiagent}. Ces propriétés correspondent au besoin de la Migration Factory, sans faire de LangGraph la source d'autorité métier.

% ============================================================
\section{Plateformes et architectures intelligentes}
% ============================================================

\subsection{Architecture en couches}

Une plateforme agentique de production ne se réduit pas à une conversation entre modèles. Elle associe généralement :

\begin{itemize}
    \item une interface utilisateur ;
    \item une API et des services applicatifs ;
    \item un orchestrateur ;
    \item des agents et des outils ;
    \item une persistance de l'état ;
    \item un magasin d'artefacts ;
    \item des mécanismes de sécurité, d'observabilité et d'évaluation.
\end{itemize}

\begin{figure}[H]
    \centering
    \resizebox{0.9\textwidth}{!}{%
    \begin{tikzpicture}[
        layer/.style={draw=cgiblue,rounded corners=3pt,fill=cgilight!55,minimum width=12.5cm,minimum height=1.0cm,align=center,font=\small\bfseries},
        arr/.style={-{Latex[length=2mm]},thick,draw=cgigray}
    ]
        \node[layer] (ui) {Interface et API};
        \node[layer,below=0.35cm of ui] (orch) {Orchestration, routage et gestion des états};
        \node[layer,below=0.35cm of orch] (agents) {Agents spécialisés et services déterministes};
        \node[layer,below=0.35cm of agents] (infra) {Persistance, artefacts, outils externes et observabilité};
        \draw[arr] (ui) -- (orch);
        \draw[arr] (orch) -- (agents);
        \draw[arr] (agents) -- (infra);
    \end{tikzpicture}%
    }
    \caption{Architecture en couches d'une plateforme agentique}
    \label{fig:plateforme-agentique-ch2}
\end{figure}

\subsection{Microservices, API et bus de messages}

Les agents peuvent être intégrés sous forme de modules internes, de services séparés ou de processus distribués. Les API définissent les contrats synchrones. Les événements et files de messages facilitent les traitements asynchrones, le découplage et la reprise. Une architecture microservices n'est cependant pas obligatoire pour obtenir un SMA : le critère essentiel est la séparation des responsabilités et des interactions, non le nombre de processus déployés.

Dans les prototypes, les agents et services sont intégrés au backend et coordonnés par le workflow. Les événements persistés jouent un rôle d'observabilité et de reconstruction. Une future industrialisation pourrait séparer certains composants ou utiliser une file de messages, mais cela n'est pas nécessaire pour le périmètre local actuel.

\subsection{Humain dans la boucle}

L'intégration d'un humain répond à trois objectifs : valider une décision à risque, apporter une information absente et arrêter une trajectoire incorrecte. Le contrôle humain ne doit pas être symbolique. Il doit porter sur un artefact précis, présenter les conséquences et enregistrer la décision.

% ============================================================
\section{Migration logicielle : motivations, types et processus}
% ============================================================

\subsection{Définitions}

\begin{table}[H]
    \centering
    \small
    \arrayrulecolor{cgigray}
    \rowcolors{2}{cgilight!55}{white}
    \begin{tabularx}{\textwidth}{|>{\bfseries\RaggedRight\arraybackslash}p{3.3cm}|>{\RaggedRight\arraybackslash}X|}
        \hline
        \rowcolor{cgilight!95}
        \textbf{Notion} & \textbf{Définition} \\
        \hline
        Mise à niveau & Passage à une version plus récente d'un composant ou d'un runtime. \\
        \hline
        Migration & Transformation contrôlée d'un système source vers une cible, incluant dépendances, configuration, code et validation. \\
        \hline
        Modernisation & Démarche plus large intégrant migration, architecture, sécurité et dette technique. \\
        \hline
        Refactorisation & Amélioration interne du code sans changement intentionnel du comportement métier. \\
        \hline
    \end{tabularx}
    \caption{Notions relatives à la transformation logicielle}
    \label{tab:notions-migration-ch2}
    \rowcolors{2}{white}{white}
    \arrayrulecolor{black}
\end{table}

\subsection{Pourquoi migrer ?}

Les migrations sont déclenchées par la fin de support, les vulnérabilités, l'incompatibilité avec les plateformes, la dette technique, les exigences de maintenabilité et l'accès à de nouvelles capacités. Le coût de l'inaction augmente lorsque les versions anciennes deviennent plus difficiles à compiler, tester ou déployer.

\subsection{Types de migration}

On distingue notamment :

\begin{itemize}
    \item migration de runtime ou de langage ;
    \item migration de framework ;
    \item mise à niveau de dépendances ;
    \item migration de plateforme ou de cloud ;
    \item migration d'architecture ;
    \item migration de données ;
    \item combinaison progressive de plusieurs transformations.
\end{itemize}

\subsection{Processus général}

\begin{figure}[H]
    \centering
    \resizebox{0.98\textwidth}{!}{%
    \begin{tikzpicture}[
        step/.style={draw=cgiblue,rounded corners=3pt,fill=cgilight!55,minimum width=2.6cm,minimum height=1.0cm,align=center,font=\small\bfseries},
        arr/.style={-{Latex[length=2mm]},thick,draw=cgigray}
    ]
        \node[step] (inventory) {Inventaire};
        \node[step,right=0.35cm of inventory] (target) {Cible};
        \node[step,right=0.35cm of target] (plan) {Plan};
        \node[step,right=0.35cm of plan] (transform) {Transformation};
        \node[step,right=0.35cm of transform] (build) {Build et tests};
        \node[step,right=0.35cm of build] (repair) {Réparation};
        \node[step,right=0.35cm of repair] (report) {Validation};
        \draw[arr] (inventory)--(target);
        \draw[arr] (target)--(plan);
        \draw[arr] (plan)--(transform);
        \draw[arr] (transform)--(build);
        \draw[arr] (build)--(repair);
        \draw[arr] (repair)--(report);
        \draw[arr,rounded corners=4pt] (repair.south)--++(0,-0.65)-|(transform.south);
    \end{tikzpicture}%
    }
    \caption{Processus général d'une migration applicative}
    \label{fig:processus-migration-ch2}
\end{figure}

\subsection{Spécificités Java/Spring Boot}

Une migration Spring Boot associe versions de Java, Spring Boot, Spring Framework, dépendances, plugins Maven ou Gradle, propriétés et espaces de noms. La documentation Spring recommande d'examiner les notes de chaque version traversée et de progresser par étapes compatibles \cite{spring-boot-upgrading}.

Le projet couvre Spring Boot 2.1 vers 2.7, 2.7 vers 3.5 et 3.5 vers 4.0, avec Java 11, 17 et 21. Les difficultés portent notamment sur les API dépréciées, la transition \texttt{javax} vers \texttt{jakarta}, les plugins, la configuration et les bibliothèques tierces.

\subsection{Spécificités Angular}

Angular dépend simultanément d'Angular CLI, Node.js, npm, TypeScript, RxJS et des bibliothèques du projet. La documentation officielle publie une matrice de compatibilité et recommande de franchir les versions majeures successivement avec l'Update Guide et \texttt{ng update} \cite{angular-versions,angular-update}.

Le périmètre couvre Angular 18 vers 19, 19 vers 20 et 20 vers 21. Chaque étape nécessite un profil Node/npm compatible, une installation propre, un build et les tests disponibles.

% ============================================================
\section{Étude de l'existant et travaux proches}
% ============================================================

\subsection{Processus manuel de référence}

Dans un processus traditionnel, le développeur analyse le projet, consulte les guides, modifie les manifestes et le code, exécute les commandes, lit les journaux puis répète le cycle. Les outils existent, mais l'état, les preuves et les décisions restent dispersés. Cette organisation limite la reproductibilité, le suivi de plusieurs migrations et la capitalisation des corrections.

\subsection{Travaux et plateformes étudiés}

Huit solutions représentatives ont été retenues afin de couvrir la migration déterministe, l'assistance intelligente, la réparation et l'orchestration multi-agents.

\paragraph{OpenRewrite.}
OpenRewrite applique des recettes composables sur des arbres sémantiques préservant le format. Les recettes peuvent réaliser des migrations de frameworks et produire des différences examinables \cite{openrewrite-recipes}. Son principal avantage est le déterminisme ; sa limite est la dépendance aux recettes disponibles.

\paragraph{Angular CLI et \texttt{ng update}.}
L'outil natif Angular applique les migrations officielles entre versions et met à jour les dépendances prises en charge. Il reste toutefois centré sur l'écosystème Angular et ne pilote pas à lui seul l'ensemble du processus de gouvernance.

\paragraph{Dependabot.}
Dependabot détecte des versions de dépendances et ouvre des pull requests pour mettre à jour les manifestes. La validation dépend ensuite des tests et de la revue de l'équipe \cite{dependabot}. Il automatise l'entretien des dépendances, non une migration complète de framework.

\paragraph{Amazon Q Developer Code Transformation.}
Amazon Q peut analyser un projet Maven, produire un plan, transformer le code Java, relancer des builds et tests, puis présenter un diff à accepter. Les scénarios supportés, les contraintes d'environnement et les limites de taille restent définis par le service \cite{amazonq-transform}.

\paragraph{GitHub Copilot Modernization.}
GitHub Copilot Modernization propose des workflows structurés d'évaluation, de planification, de transformation et de validation pour Java et d'autres environnements. Les fonctions récentes prennent en charge plusieurs scénarios Spring Boot et des workflows multi-agents \cite{copilot-modernization}.

\paragraph{SWE-agent.}
SWE-agent fournit à un agent LLM une interface spécialisée pour naviguer dans un dépôt, modifier des fichiers et exécuter des tests. Son évaluation sur SWE-bench montre l'importance de l'interface agent--ordinateur, mais son objectif est la résolution autonome d'issues plutôt que la gouvernance d'une trajectoire de migration \cite{swe-agent}.

\paragraph{AutoGen.}
AutoGen organise des conversations entre agents configurables combinant LLM, outils et interventions humaines. Le framework illustre plusieurs patrons de collaboration et de review, mais laisse au concepteur la responsabilité de la persistance, de la sécurité et de l'intégration métier \cite{autogen-paper}.

\paragraph{LangGraph.}
LangGraph fournit un graphe d'exécution avec état, routage, checkpoints et interruptions. Il est adapté aux workflows hybrides, mais ne définit pas les règles de migration, les agents ou la gouvernance ; ces éléments doivent être conçus par l'application.

\subsection{Comparaison synthétique}

\scriptsize
\begin{longtable}{|>{\bfseries\RaggedRight\arraybackslash}p{2.8cm}|>{\RaggedRight\arraybackslash}p{2.6cm}|>{\RaggedRight\arraybackslash}p{2.5cm}|>{\RaggedRight\arraybackslash}p{2.5cm}|>{\RaggedRight\arraybackslash}p{2.5cm}|>{\RaggedRight\arraybackslash}p{3.2cm}|}
    \caption{Comparaison des travaux et plateformes proches}\label{tab:travaux-proches-ch2}\\
    \hline
    \rowcolor{cgilight!95}
    \textbf{Travail} & \textbf{Domaine} & \textbf{Agents} & \textbf{Orchestration} & \textbf{Évaluation} & \textbf{Limite principale} \\
    \hline
    \endfirsthead
    \hline
    \rowcolor{cgilight!95}
    \textbf{Travail} & \textbf{Domaine} & \textbf{Agents} & \textbf{Orchestration} & \textbf{Évaluation} & \textbf{Limite principale} \\
    \hline
    \endhead
    OpenRewrite & Refactorisation et migration & Aucun LLM requis & Pipeline de recettes & Diff et tests du projet & Cas limités au catalogue de recettes \\
    \hline
    Angular CLI & Mise à niveau Angular & Aucun agent explicite & Séquence officielle de migrations & Build et tests externes & Périmètre Angular uniquement \\
    \hline
    Dependabot & Dépendances & Bot spécialisé & Planification périodique & Revue de PR et CI & Pas de migration applicative complète \\
    \hline
    Amazon Q & Transformation Java & Agent de transformation & Service géré, plan et itérations & Build, tests et diff & Scénarios et quotas supportés \\
    \hline
    Copilot Modernization & Modernisation applicative & Agents spécialisés & Workflow agentique structuré & Build, tests, plan et rapport & Service propriétaire et périmètre supporté \\
    \hline
    SWE-agent & Résolution d'issues & Agent d'exécution & Boucle autonome avec outils & SWE-bench, pass@1 & Gouvernance métier limitée \\
    \hline
    AutoGen & Applications multi-agents & Agents conversationnels & Conversations et patterns configurables & Benchmarks et études de cas & Infrastructure métier à construire \\
    \hline
    LangGraph & Orchestration agentique & Définis par l'application & Graphe d'état et checkpoints & Traces et tests applicatifs & Ne fournit pas les agents métier \\
    \hline
    Migration Factory & Migration Java et Angular & Agents spécialisés et reviewers & Orchestration centrale hybride & Build, tests, preuves, décisions et rapports & Prototype local à industrialiser \\
    \hline
\end{longtable}
\normalsize

\subsection{Carte de positionnement}

\begin{figure}[H]
    \centering
    \resizebox{0.93\textwidth}{!}{%
    \begin{tikzpicture}[
        root/.style={draw=cgired,rounded corners=3pt,fill=white,minimum width=4.1cm,minimum height=1.0cm,align=center,font=\small\bfseries},
        cat/.style={draw=cgiblue,rounded corners=3pt,fill=cgilight!55,minimum width=4.1cm,minimum height=1.0cm,align=center,font=\small\bfseries},
        item/.style={draw=cgigray,rounded corners=3pt,fill=white,text width=3.7cm,minimum height=1.2cm,align=center,font=\scriptsize},
        arr/.style={-{Latex[length=2mm]},thick,draw=cgigray}
    ]
        \node[root] (root) {Écosystème étudié};
        \node[cat,below left=1.2cm and 3.8cm of root] (det) {Migration déterministe};
        \node[cat,below=1.2cm of root] (ai) {Assistance intelligente};
        \node[cat,below right=1.2cm and 3.8cm of root] (mas) {Orchestration multi-agents};
        \node[item,below=0.55cm of det] {OpenRewrite\\Angular CLI\\Dependabot};
        \node[item,below=0.55cm of ai] {Amazon Q\\Copilot Modernization\\SWE-agent};
        \node[item,below=0.55cm of mas] {AutoGen\\LangGraph\\FIPA};
        \draw[arr] (root)--(det);
        \draw[arr] (root)--(ai);
        \draw[arr] (root)--(mas);
    \end{tikzpicture}%
    }
    \caption{Regroupement des approches étudiées}
    \label{fig:mindmap-etat-art-ch2}
\end{figure}

% ============================================================
\section{Positionnement de la Migration Factory}
% ============================================================

La Migration Factory se situe à l'intersection des trois familles. Elle utilise les transformations déterministes lorsque les règles sont connues, les agents lorsque le problème demande une analyse contextuelle, et un orchestrateur pour maintenir l'ordre, l'état et les décisions.

Sa contribution principale n'est pas de proposer un nouveau modèle de langage ni de remplacer les outils existants. Elle consiste à les organiser dans une architecture multi-agents gouvernée qui :

\begin{itemize}
    \item spécialise les agents par responsabilité ;
    \item sépare proposition, review, décision et exécution ;
    \item associe les sorties aux preuves et empreintes ;
    \item protège la source par des espaces isolés ;
    \item valide par les commandes réelles de l'écosystème ;
    \item conserve l'état et permet la reprise ;
    \item fournit une vue de pilotage au développeur.
\end{itemize}

\section*{Conclusion}
\addcontentsline{toc}{section}{Conclusion}

Ce chapitre a présenté les notions nécessaires à la compréhension d'une architecture multi-agents appliquée au génie logiciel. Un agent combine un objectif, un contexte, des capacités et des règles d'interaction. Un SMA coordonne plusieurs agents spécialisés, mais introduit aussi des coûts de communication, d'état et de gouvernance. Le modèle retenu par le projet est hybride et centralement orchestré : les agents contribuent à l'analyse et à la proposition, tandis que les services déterministes et le développeur conservent l'autorité.

L'étude de la migration et des travaux proches montre qu'aucune solution examinée ne réunit à elle seule spécialisation multi-agents, trajectoires Java et Angular, isolation, preuves, réparation gouvernée et supervision humaine. Le chapitre suivant traduit ce positionnement en une conception détaillée de l'architecture SMA, de ses agents, de son orchestrateur et de ses interactions.
